%!TEX TS-program = pdflatex
\documentclass[aspectratio=169,10pt]{beamer}

\usepackage[T2A]{fontenc}
\usepackage[utf8]{inputenc}
\usepackage[english,russian]{babel}
\usepackage{cmap}
\usepackage{amsmath,amssymb,mathtools}
\usepackage{booktabs}
\usepackage{array}
\usepackage{etoolbox}
\usepackage{PTSans}
\usepackage{graphicx}
\graphicspath{{figures/}}

\renewcommand{\familydefault}{\sfdefault}
\renewcommand{\rmdefault}{\sfdefault}
\newbool{russian}
\booltrue{russian}
\usepackage{HSE-theme/beamerthemeHSE}

\hypersetup{bookmarks=false}
\setbeamertemplate{blocks}[rounded][shadow=false]
\setbeamerfont{frametitle}{size=\Large}
\setbeamerfont{framesubtitle}{size=\small}
\setbeamerfont{block title}{size=\normalsize}
\setbeamerfont{block body}{size=\small}

\newcommand{\R}{\mathbb{R}}
\newcommand{\norm}[1]{\left\lVert #1\right\rVert}
\newcommand{\grad}{\nabla}
\newcommand{\gap}[1]{f(#1)-f^\star}

\title[Методы оптимизации в ML]{Методы оптимизации\\в машинном обучении}
\subtitle{Лекция 6. Скорость сходимости градиентного спуска}
\author[Н. Лукьяненко]{Никита Лукьяненко}
\institute{Совместная программа СМОЛ ГУ и ФКН НИУ ВШЭ}
\date{2026/27 учебный год}

\begin{document}

% 1
\frame[plain]{\titlepage}

% 2
\begin{frame}{Что мы уже знаем про градиентный спуск}
\small
Для $L$-гладкой функции при шаге $\alpha=1/L$:
\[
 x_{k+1}=x_k-\frac1L\grad f(x_k)
\]
и
\[
\boxed{
 f(x_{k+1})
 \le
 f(x_k)-\frac{1}{2L}\norm{\grad f(x_k)}^2.
}
\]
Следовательно, значения функции не возрастают.

\begin{alertblock}{Но этого мало}
Монотонное убывание ещё не говорит, \textbf{насколько быстро} мы приближаемся к оптимуму.
\end{alertblock}
\end{frame}

% 3
\begin{frame}{Из лекции 5: теперь мы знаем, куда хотим прийти}
\small
Для дифференцируемой выпуклой функции
\[
\grad f(x^\star)=0
\quad\Longrightarrow\quad
x^\star\text{ — глобальный минимум.}
\]
Для сильно выпуклой функции этот минимум единственен.

\vspace{0.8em}
Сегодня вопрос уже не в существовании хорошего решения, а в сложности его поиска:
\[
\boxed{\text{Сколько итераций нужно градиентному спуску?}}
\]
\end{frame}

% 4
\begin{frame}{Три разные меры ошибки}
\small
На итерации $k$ можно измерять качество по-разному:
\[
\boxed{f(x_k)-f^\star}
\qquad
\boxed{\norm{x_k-x^\star}}
\qquad
\boxed{\norm{\grad f(x_k)}}.
\]

\centering
\includegraphics[width=0.80\linewidth]{error_metrics.pdf}
\end{frame}

% 5
\begin{frame}{Они связаны, но не совпадают}
\small
Одинаковое расстояние до $x^\star$ может давать совершенно разную ошибку по функции — всё зависит от кривизны.

А маленький градиент сам по себе без дополнительных предположений не обязан означать маленькую ошибку по функции.

\begin{block}{Основная мера сегодня}
Будем оценивать
\[
\boxed{f(x_k)-f^\star.}
\]
Это непосредственно отвечает на вопрос, насколько хорошее значение целевой функции уже найдено.
\end{block}
\end{frame}

% 6
\begin{frame}{Два режима, которые мы получим}
\centering
\includegraphics[width=0.86\linewidth]{rate_preview.pdf}

\vspace{0.2em}
\small
Для выпуклой гладкой функции появится скорость порядка $1/k$, а сильная выпуклость даст геометрическое уменьшение ошибки.
\end{frame}

% 7
\begin{frame}{Главные результаты лекции — пока без доказательства}
\small
Для выпуклой $L$-гладкой функции:
\[
\boxed{
 f(x_k)-f^\star
 \le
 \frac{L\norm{x_0-x^\star}^2}{2k}.
}
\]
Для $\mu$-сильно выпуклой и $L$-гладкой:
\[
\boxed{
 f(x_k)-f^\star
 \le
 \left(1-\frac{\mu}{L}\right)^k
 \bigl(f(x_0)-f^\star\bigr).
}
\]
Наша задача — понять, откуда возникают обе формулы.
\end{frame}

% 8
\begin{frame}{$L$-гладкость: ограничение на максимальную кривизну}
\small
Напомним определение:
\[
\norm{\grad f(x)-\grad f(y)}
\le
L\norm{x-y}.
\]
Если $f$ дважды дифференцируема, то интуитивно
\[
L\approx\text{максимальная кривизна функции}.
\]
Для квадратичной функции
\[
f(x)=\frac12x^\top Ax-b^\top x
\]
имеем
\[
L=\lambda_{\max}(A).
\]
\end{frame}

% 9
\begin{frame}{Главный инструмент из лекции 4}
\small
Из $L$-гладкости следует квадратичная верхняя оценка:
\[
f(y)
\le
f(x)+\grad f(x)^\top(y-x)
+\frac L2\norm{y-x}^2.
\]
При градиентном шаге $y=x-\frac1L\grad f(x)$:
\[
\boxed{
 f(x_{k+1})
 \le
 f(x_k)-\frac1{2L}\norm{\grad f(x_k)}^2.
}
\]
\end{frame}

% 10
\begin{frame}{Главный инструмент из лекции 5}
\small
Для выпуклой дифференцируемой функции
\[
 f(y)
 \ge
 f(x)+\grad f(x)^\top(y-x).
\]
Это означает: касательная в $x$ — глобальная нижняя оценка функции.

\begin{block}{Нам понадобится одна специальная подстановка}
Положим $y=x^\star$.
\end{block}
\end{frame}

% 11
\begin{frame}{Выпуклость связывает градиент с ошибкой по функции}
\small
Подставляя $y=x^\star$, получаем
\[
f^\star
\ge
f(x)+\grad f(x)^\top(x^\star-x).
\]
После переноса:
\[
\boxed{
 f(x)-f^\star
 \le
 \grad f(x)^\top(x-x^\star).
}
\]
Справа появилась информация о направлении на оптимум.
\end{frame}

% 12
\begin{frame}{Сильная выпуклость добавляет квадратичный запас}
\small
Если $f$ $\mu$-сильно выпукла, то
\[
f(y)
\ge
f(x)+\grad f(x)^\top(y-x)
+\frac\mu2\norm{y-x}^2.
\]
Это более сильное утверждение, чем обычная выпуклость.

\centering
\includegraphics[width=0.62\linewidth]{strong_lower_model.pdf}
\end{frame}

% 13
\begin{frame}{Карта доказательств}
\small
В первой части:
\[
\boxed{L\text{-гладкость}+\text{выпуклость}}
\Longrightarrow
\boxed{f(x_k)-f^\star=O(1/k)}.
\]
Во второй части добавим сильную выпуклость:
\[
\boxed{L\text{-гладкость}+\mu\text{-сильная выпуклость}}
\Longrightarrow
\boxed{f(x_k)-f^\star\le q^k C}.
\]
И увидим, что
\[
q=1-\frac{1}{\kappa},
\qquad
\kappa=\frac{L}{\mu}.
\]
\end{frame}

% 14
\begin{frame}{Градиентный шаг даёт полезное равенство}
\small
Работаем с шагом
\[
x_{k+1}=x_k-\frac1L\grad f(x_k).
\]
Переносим:
\[
\frac1L\grad f(x_k)=x_k-x_{k+1}.
\]
Поэтому
\[
\boxed{
\grad f(x_k)=L(x_k-x_{k+1}).
}
\]
Эта простая формула позволит заменить градиент геометрией последовательных точек.
\end{frame}

% 15
\begin{frame}{Какой вид оценки нам хотелось бы получить?}
\small
Если удастся доказать неравенство вида
\[
\boxed{
 f(x_{k+1})-f^\star
 \le
 C\Bigl(
 \norm{x_k-x^\star}^2
 -
 \norm{x_{k+1}-x^\star}^2
 \Bigr),
}
\]
то при сложении по $k$ большая часть членов сократится.

\begin{alertblock}{Именно это и произойдёт}
Такие разности называются телескопическими: внутренние элементы исчезают при суммировании.
\end{alertblock}
\end{frame}

% 16
\begin{frame}{Шаг 1: применяем гладкость}
\small
К точкам $x_k$ и $x_{k+1}$:
\[
\begin{aligned}
f(x_{k+1})
\le{}&
f(x_k)
+\grad f(x_k)^\top(x_{k+1}-x_k)\\
&+\frac L2\norm{x_{k+1}-x_k}^2.
\end{aligned}
\]
Вычитаем $f^\star$:
\[
\gap{x_{k+1}}
\le
\gap{x_k}
+\grad f(x_k)^\top(x_{k+1}-x_k)
+\frac L2\norm{x_{k+1}-x_k}^2.
\]
\end{frame}

% 17
\begin{frame}{Шаг 2: используем выпуклость}
\small
Из слайда 11:
\[
\gap{x_k}
\le
\grad f(x_k)^\top(x_k-x^\star).
\]
Подставляем:
\[
\begin{aligned}
\gap{x_{k+1}}
\le{}&
\grad f(x_k)^\top(x_k-x^\star)\\
&+\grad f(x_k)^\top(x_{k+1}-x_k)\\
&+\frac L2\norm{x_{k+1}-x_k}^2.
\end{aligned}
\]
\end{frame}

% 18
\begin{frame}{Два скалярных произведения объединяются}
\small
Первые два члена имеют общий градиент:
\[
\begin{aligned}
&\grad f(x_k)^\top(x_k-x^\star)
+\grad f(x_k)^\top(x_{k+1}-x_k)\\
&\qquad=
\grad f(x_k)^\top(x_{k+1}-x^\star).
\end{aligned}
\]
Поэтому
\[
\boxed{
\gap{x_{k+1}}
\le
\grad f(x_k)^\top(x_{k+1}-x^\star)
+\frac L2\norm{x_{k+1}-x_k}^2.
}
\]
\end{frame}

% 19
\begin{frame}{Шаг 3: заменяем градиент через шаг}
\small
Используем
\[
\grad f(x_k)=L(x_k-x_{k+1}).
\]
Получаем
\[
\begin{aligned}
\gap{x_{k+1}}
\le{}&
L(x_k-x_{k+1})^\top(x_{k+1}-x^\star)\\
&+\frac L2\norm{x_k-x_{k+1}}^2.
\end{aligned}
\]
Осталось распознать чисто геометрическое тождество.
\end{frame}

% 20
\begin{frame}{Трёхточечное тождество}
\small
Для любых $a,b,c$:
\[
\boxed{
2(a-b)^\top(b-c)
=
\norm{a-c}^2
-\norm{a-b}^2
-\norm{b-c}^2.
}
\]
В нашем случае:
\[
a=x_k,
\qquad
b=x_{k+1},
\qquad
c=x^\star.
\]
\end{frame}

% 21
\begin{frame}{Почему тождество верно}
\small
Заметим:
\[
a-c=(a-b)+(b-c).
\]
Тогда
\[
\begin{aligned}
\norm{a-c}^2
&=\norm{a-b}^2+\norm{b-c}^2\\
&\quad+2(a-b)^\top(b-c).
\end{aligned}
\]
Перенос двух первых слагаемых даёт нужную формулу.

\begin{block}{Никакой новой теории}
Это просто раскрытие квадрата нормы суммы двух векторов.
\end{block}
\end{frame}

% 22
\begin{frame}{Ключевое неравенство для выпуклого случая}
\small
Подставляя трёхточечное тождество, получаем
\[
\begin{aligned}
\gap{x_{k+1}}
\le
\frac L2\Bigl(&
\norm{x_k-x^\star}^2
-\norm{x_{k+1}-x^\star}^2\\
&-\norm{x_k-x_{k+1}}^2\Bigr)
+\frac L2\norm{x_k-x_{k+1}}^2.
\end{aligned}
\]
Последние два члена сокращаются:
\[
\boxed{
\gap{x_{k+1}}
\le
\frac L2
\left(
\norm{x_k-x^\star}^2-
\norm{x_{k+1}-x^\star}^2
\right).
}
\]
\end{frame}

% 23
\begin{frame}{Теперь складываем неравенства}
\small
Для $k=0,\ldots,N-1$:
\[
\sum_{k=0}^{N-1}
\bigl(f(x_{k+1})-f^\star\bigr)
\le
\frac L2
\sum_{k=0}^{N-1}
\left(
\norm{x_k-x^\star}^2-
\norm{x_{k+1}-x^\star}^2
\right).
\]
Справа возникает телескопическая сумма.
\end{frame}

% 24
\begin{frame}{Телескопическая сумма}
\centering
\includegraphics[width=0.88\linewidth]{telescoping_sum.pdf}

\vspace{0.2em}
\small
В результате остаются только начальное и конечное расстояния.
\end{frame}

% 25
\begin{frame}{После суммирования}
\small
Получаем
\[
\sum_{k=1}^{N}
\bigl(f(x_k)-f^\star\bigr)
\le
\frac L2
\left(
\norm{x_0-x^\star}^2-
\norm{x_N-x^\star}^2
\right).
\]
Поскольку норма неотрицательна,
\[
\boxed{
\sum_{k=1}^{N}
\bigl(f(x_k)-f^\star\bigr)
\le
\frac L2
\norm{x_0-x^\star}^2.
}
\]
\end{frame}

% 26
\begin{frame}{Используем монотонность значений функции}
\small
По лемме об убывании
\[
f(x_N)
\le
f(x_{N-1})
\le\cdots\le
f(x_1).
\]
Значит каждое из первых $N$ отклонений не меньше последнего:
\[
 f(x_k)-f^\star
 \ge
 f(x_N)-f^\star.
\]
Поэтому
\[
\boxed{
N\bigl(f(x_N)-f^\star\bigr)
\le
\sum_{k=1}^{N}
\bigl(f(x_k)-f^\star\bigr).
}
\]
\end{frame}

% 27
\begin{frame}{Теорема: скорость GD для выпуклой гладкой функции}
\small
\begin{block}{Теорема}
Пусть $f$ выпукла и $L$-гладка, а
\[
x_{k+1}=x_k-\frac1L\grad f(x_k).
\]
Тогда для любого $N\ge1$
\[
\boxed{
 f(x_N)-f^\star
 \le
 \frac{L\norm{x_0-x^\star}^2}{2N}.
}
\]
\end{block}
Это и есть скорость порядка $O(1/N)$.
\end{frame}

% 28
\begin{frame}{Доказательство теоремы — в трёх строках}
\footnotesize
Из ключевого неравенства:
\[
\sum_{k=1}^{N}
\bigl(f(x_k)-f^\star\bigr)
\le
\frac L2\norm{x_0-x^\star}^2.
\]
Из монотонности:
\[
N\bigl(f(x_N)-f^\star\bigr)
\le
\sum_{k=1}^{N}
\bigl(f(x_k)-f^\star\bigr).
\]
Следовательно,
\[
N\bigl(f(x_N)-f^\star\bigr)
\le
\frac L2\norm{x_0-x^\star}^2,
\]
откуда делением на $N$ получается утверждение.
\end{frame}

% 29
\begin{frame}{Что означает $O(1/k)$?}
\centering
\includegraphics[width=0.70\linewidth]{one_over_k.pdf}

\vspace{0.2em}
\small
При удвоении числа итераций гарантированная верхняя оценка примерно уменьшается вдвое.
\end{frame}

% 30
\begin{frame}{Сколько итераций нужно для точности $\varepsilon$?}
\small
Хотим обеспечить
\[
f(x_k)-f^\star\le\varepsilon.
\]
Из теоремы достаточно потребовать
\[
\frac{L\norm{x_0-x^\star}^2}{2k}
\le
\varepsilon.
\]
Следовательно,
\[
\boxed{
k
\ge
\frac{L\norm{x_0-x^\star}^2}{2\varepsilon}.
}
\]
То есть сложность по точности — порядка $O(1/\varepsilon)$.
\end{frame}

% 31
\begin{frame}{Почему такая скорость может быть медленной}
\small
Если хотим улучшить точность в десять раз,
\[
\varepsilon
\longrightarrow
\frac{\varepsilon}{10},
\]
то верхняя оценка числа итераций увеличивается примерно в десять раз.

\begin{exampleblock}{Пример}
Переход от ошибки $10^{-2}$ к $10^{-3}$ в худшей гарантии требует на порядок большего $k$.
\end{exampleblock}
\end{frame}

% 32
\begin{frame}{Важно: это верхняя оценка, а не точный прогноз}
\small
Теорема утверждает
\[
f(x_k)-f^\star
\le
\frac{C}{k}.
\]
Она \textbf{не} утверждает, что
\[
f(x_k)-f^\star=\frac{C}{k}.
\]

\begin{alertblock}{Как читать оценки сходимости}
Они гарантируют: алгоритм работает \textbf{не хуже определённого сценария} на всём рассматриваемом классе функций. На конкретной задаче он может быть гораздо быстрее.
\end{alertblock}
\end{frame}

% 33
\begin{frame}{Теперь добавим сильную выпуклость}
\small
Пусть $f$ одновременно
\begin{itemize}
\item $L$-гладкая;
\item $\mu$-сильно выпуклая.
\end{itemize}
Тогда
\[
0<\mu\le L.
\]
Сильная выпуклость даст связь между величиной градиента и ошибкой по функции — именно её не хватало раньше.
\end{frame}

% 34
\begin{frame}{Что именно даёт сильная выпуклость}
\small
Для любого $y$:
\[
f(y)
\ge
f(x)+\grad f(x)^\top(y-x)
+\frac\mu2\norm{y-x}^2.
\]
Правая часть — квадратичная функция по $y$.

\begin{block}{Идея}
Если она является глобальной нижней оценкой $f$, можно минимизировать её и получить нижнюю оценку на само $f^\star$.
\end{block}
\end{frame}

% 35
\begin{frame}{Минимизируем квадратичную нижнюю модель}
\small
Положим
\[
h=y-x.
\]
Нужно минимизировать по $h$
\[
q(h)=\grad f(x)^\top h+\frac\mu2\norm{h}^2.
\]
Условие стационарности:
\[
\grad_h q(h)=\grad f(x)+\mu h=0.
\]
Следовательно,
\[
\boxed{
h^\star=-\frac1\mu\grad f(x).
}
\]
\end{frame}

% 36
\begin{frame}{Минимальное значение модели}
\small
Подставляем $h^\star=-\frac1\mu\grad f(x)$:
\[
\begin{aligned}
q(h^\star)
&=-\frac1\mu\norm{\grad f(x)}^2
+\frac\mu2\frac1{\mu^2}\norm{\grad f(x)}^2\\
&=-\frac1{2\mu}\norm{\grad f(x)}^2.
\end{aligned}
\]
Поэтому для любого $y$ и, в частности, для минимума:
\[
f^\star
\ge
f(x)-\frac1{2\mu}\norm{\grad f(x)}^2.
\]
\end{frame}

% 37
\begin{frame}{Сильная выпуклость не позволяет градиенту быть слишком маленьким}
\small
Переносим члены:
\[
 f(x)-f^\star
 \le
 \frac1{2\mu}\norm{\grad f(x)}^2.
\]
Или эквивалентно
\[
\boxed{
\norm{\grad f(x)}^2
\ge
2\mu\bigl(f(x)-f^\star\bigr).
}
\]
Это ключевой мост между леммой об убывании и ошибкой по функции.
\end{frame}

% 38
\begin{frame}{Геометрический смысл предыдущего неравенства}
\centering
\includegraphics[width=0.82\linewidth]{gradient_gap_relation.pdf}

\vspace{0.15em}
\small
В сильно выпуклой задаче заметное отставание по функции обязано сопровождаться достаточно большим градиентом.
\end{frame}

% 39
\begin{frame}{Снова используем лемму об убывании}
\small
При шаге $1/L$:
\[
\gap{x_{k+1}}
\le
\gap{x_k}
-
\frac1{2L}\norm{\grad f(x_k)}^2.
\]
Из сильной выпуклости:
\[
\norm{\grad f(x_k)}^2
\ge
2\mu\bigl(f(x_k)-f^\star\bigr).
\]
Подставляем нижнюю оценку на градиент в отрицательный член.
\end{frame}

% 40
\begin{frame}{Получаем сокращение ошибки за один шаг}
\small
\[
\begin{aligned}
\gap{x_{k+1}}
&\le
\gap{x_k}
-
\frac1{2L}\cdot2\mu\bigl(f(x_k)-f^\star\bigr)\\
&=
\left(1-\frac\mu L\right)\bigl(f(x_k)-f^\star\bigr).
\end{aligned}
\]
То есть
\[
\boxed{
\gap{x_{k+1}}
\le
\left(1-\frac\mu L\right)\bigl(f(x_k)-f^\star\bigr).
}
\]
Каждая итерация сокращает верхнюю оценку на один и тот же множитель.
\end{frame}

% 41
\begin{frame}{Повторяем неравенство $k$ раз}
\small
Последовательно:
\[
\gap{x_1}
\le
q\bigl(f(x_0)-f^\star\bigr),
\]
\[
\gap{x_2}
\le
q\bigl(f(x_1)-f^\star\bigr)
\le
q^2\bigl(f(x_0)-f^\star\bigr),
\]
где
\[
q=1-\frac\mu L.
\]
Поэтому
\[
\boxed{
\gap{x_k}
\le
\left(1-\frac\mu L\right)^k\bigl(f(x_0)-f^\star\bigr).
}
\]
\end{frame}

% 42
\begin{frame}{Теорема о линейной сходимости}
\small
\begin{block}{Теорема}
Пусть $f$ $\mu$-сильно выпукла и $L$-гладка. Для GD с шагом $1/L$:
\[
\boxed{
 f(x_k)-f^\star
 \le
 \left(1-\frac\mu L\right)^k
 \bigl(f(x_0)-f^\star\bigr).
}
\]
\end{block}

\begin{alertblock}{Терминология}
«Линейная сходимость» не означает движение по прямой. Это стандартное название геометрического закона $e_{k+1}\le q e_k$, где $q<1$.
\end{alertblock}
\end{frame}

% 43
\begin{frame}{$1/k$ против геометрической скорости}
\centering
\includegraphics[width=0.88\linewidth]{rate_preview.pdf}

\vspace{0.15em}
\small
Геометрическое уменьшение особенно хорошо видно на логарифмической шкале: оно становится почти прямой линией.
\end{frame}

% 44
\begin{frame}{Число обусловленности снова появляется}
\small
Определим
\[
\boxed{\kappa=\frac{L}{\mu}.}
\]
Тогда множитель сокращения
\[
q
=1-\frac\mu L
=1-\frac1\kappa.
\]
Следовательно,
\[
\boxed{
\gap{x_k}
\le
\left(1-\frac1\kappa\right)^k\bigl(f(x_0)-f^\star\bigr).
}
\]
\end{frame}

% 45
\begin{frame}{Хорошо обусловленная задача}
\small
Пусть
\[
\kappa=2.
\]
Тогда
\[
q=1-\frac12=0.5.
\]
То есть гарантированная верхняя оценка ошибки делится примерно пополам на каждой итерации:
\[
1,\;0.5,\;0.25,\;0.125,\ldots
\]
\end{frame}

% 46
\begin{frame}{Плохо обусловленная задача}
\small
Пусть
\[
\kappa=100.
\]
Тогда
\[
q=1-\frac1{100}=0.99.
\]
За одну итерацию гарантированно исчезает лишь около $1\%$ текущей ошибки.

\centering
\includegraphics[width=0.62\linewidth]{kappa_rates.pdf}
\end{frame}

% 47
\begin{frame}{Сколько итераций нужно в сильно выпуклом случае?}
\small
Хотим
\[
\left(1-\frac1\kappa\right)^k
\le
\varepsilon.
\]
Используем
\[
1-u\le e^{-u},
\qquad u\ge0.
\]
Тогда
\[
\left(1-\frac1\kappa\right)^k
\le
e^{-k/\kappa}.
\]
\end{frame}

% 48
\begin{frame}{Откуда возникает логарифм}
\small
Чтобы
\[
e^{-k/\kappa}\le\varepsilon,
\]
достаточно
\[
-\frac{k}{\kappa}\le\log\varepsilon.
\]
То есть
\[
\boxed{
k\ge\kappa\log\frac1\varepsilon.}
\]
В обозначениях порядка:
\[
\boxed{
k=O\!\left(\kappa\log\frac1\varepsilon\right).}
\]
\end{frame}

% 49
\begin{frame}{Сравнение двух классов функций}
\small
\centering
\begin{tabular}{lcc}
\toprule
& Выпуклая + гладкая & Сильно выпуклая + гладкая\\
\midrule
Ошибка & $O(1/k)$ & $O(q^k)$\\
До точности $\varepsilon$ & $O(1/\varepsilon)$ & $O(\log(1/\varepsilon))$\\
Минимум единственен & не обязательно & да\\
Роль $\kappa$ & не обязательна & определяет скорость\\
\bottomrule
\end{tabular}

\vspace{0.7em}
\includegraphics[width=0.53\linewidth]{complexity_epsilon.pdf}
\end{frame}

% 50
\begin{frame}{Связь с квадратичными задачами}
\small
Для
\[
f(x)=\frac12x^\top Ax-b^\top x,
\qquad A\succ0,
\]
имеем
\[
L=\lambda_{\max}(A),
\qquad
\mu=\lambda_{\min}(A).
\]
Поэтому
\[
\boxed{
\kappa=\frac{L}{\mu}
=\frac{\lambda_{\max}(A)}{\lambda_{\min}(A)}.
}
\]
Это ровно то же число обусловленности, которое появилось на лекции 3.
\end{frame}

% 51
\begin{frame}{Один минимум, разные числа обусловленности}
\centering
\includegraphics[width=0.96\linewidth]{gd_trajectories_kappa.pdf}

\vspace{0.1em}
\small
Чем более вытянуты линии уровня, тем дольше градиентный спуск избавляется от медленной компоненты ошибки.
\end{frame}

% 52
\begin{frame}{Фактическая ошибка и теоретическая граница}
\centering
\includegraphics[width=0.76\linewidth]{objective_gaps.pdf}

\vspace{0.15em}
\small
Пунктир — гарантированная скорость $\left(1-1/\kappa\right)^k$; сплошные линии — реальный эксперимент.
\end{frame}

% 53
\begin{frame}{Зачем используют логарифмический масштаб}
\small
Если
\[
e_k=Cq^k,
\qquad 0<q<1,
\]
то после логарифмирования
\[
\log e_k=\log C+k\log q.
\]
Это линейная функция от $k$.

\begin{block}{Практический признак}
Почти прямая линия ошибки на полулогарифмическом графике — характерный признак геометрической сходимости.
\end{block}
\end{frame}

% 54
\begin{frame}{Теория против эксперимента}
\centering
\includegraphics[width=0.73\linewidth]{theory_vs_experiment.pdf}

\vspace{0.15em}
\small
Теоретическая оценка строится для всего класса задач и потому часто консервативна. Она должна быть над фактической ошибкой, а не совпадать с ней.
\end{frame}

% 55
\begin{frame}{Самопроверка}
\small
Пусть
\[
L=20,
\qquad
\mu=2.
\]
Ответьте:
\begin{enumerate}
\item Чему равно $\kappa$?
\item Чему равен множитель $q=1-1/\kappa$?
\item Во сколько раз примерно уменьшится верхняя оценка ошибки после 20 шагов?
\end{enumerate}

\[
\boxed{\kappa=10,\qquad q=0.9,\qquad q^{20}\approx0.12.}
\]
\end{frame}

% 56
\begin{frame}{Карта лекции}
\small
Для выпуклой $L$-гладкой функции:
\[
\boxed{
\gap{x_k}
\le
\frac{L\norm{x_0-x^\star}^2}{2k}.
}
\]
Добавляем $\mu$-сильную выпуклость:
\[
\boxed{
\gap{x_k}
\le
\left(1-\frac\mu L\right)^k\bigl(f(x_0)-f^\star\bigr).
}
\]
Именно отношение
\[
\boxed{\kappa=L/\mu}
\]
управляет худшей гарантированной геометрической скоростью.
\end{frame}

% 57
\begin{frame}{Следующий естественный вопрос}
\small
При плохой обусловленности
\[
\kappa\gg1
\]
обычный градиентный спуск может быть медленным, даже если задача выпуклая и очень хорошо изучена.

\vspace{0.6em}
\begin{alertblock}{Что можно сделать?}
Можно ли изменить алгоритм так, чтобы быстрее проходить длинные узкие долины и меньше зависеть от $\kappa$?
\end{alertblock}

Это приведёт нас к более продвинутым методам оптимизации.
\end{frame}


\end{document}
